\chapter{推导、可靠性、完备性与紧致性}
\begin{qlnote}{阅读参照}Enderton, \emph{Logic}, Chapters One--Two。
下面用自然演绎展示推理；阅读原书时另行对照其公理式演算。\end{qlnote}

\section{有限推导与假设的解除}
\(\Gamma\vdash\varphi\) 表示存在一个有限的形式推导，从 \(\Gamma\) 中的假设推出
\(\varphi\)。推导的每一步须由规定的规则允许；推导树的叶子记录尚未解除的假设。
以下规则中的 \(\neg\varphi\) 可理解为 \(\varphi\to\bot\)。

\begin{align*}
\frac{\Gamma\vdash\varphi\quad\Gamma\vdash\psi}{\Gamma\vdash\varphi\land\psi}
&\qquad
\frac{\Gamma\vdash\varphi\land\psi}{\Gamma\vdash\varphi}
\qquad
\frac{\Gamma\vdash\varphi\land\psi}{\Gamma\vdash\psi},\\
\frac{\Gamma,\varphi\vdash\psi}{\Gamma\vdash\varphi\to\psi}
&\qquad
\frac{\Gamma\vdash\varphi\to\psi\quad\Gamma\vdash\varphi}{\Gamma\vdash\psi},\\
\frac{\Gamma\vdash\varphi}{\Gamma\vdash\varphi\lor\psi}
&\qquad
\frac{\Gamma\vdash\psi}{\Gamma\vdash\varphi\lor\psi},\\
\frac{\Gamma\vdash\varphi\lor\psi\quad\Gamma,\varphi\vdash\chi\quad\Gamma,\psi\vdash\chi}
 {\Gamma\vdash\chi}
&\qquad
\frac{\Gamma\vdash\bot}{\Gamma\vdash\chi}.
\end{align*}
假设规则允许从 \(\Gamma,\varphi\) 推出 \(\varphi\)；背景假设可加入而不必使用。
条件引入和分情况证明会解除相应的临时假设。
经典反证规则还允许从 \(\Gamma,\neg\varphi\vdash\bot\) 推出 \(\Gamma\vdash\varphi\)。

\section{量词规则的附加条件}
全称消去允许由 \(\forall x\varphi\) 推出 \(\varphi[t/x]\)，前提是代入不捕获变量。
全称引入允许由 \(\Gamma\vdash\varphi\) 推出 \(\Gamma\vdash\forall x\varphi\)，
但 \(x\) 不能在 \(\Gamma\) 的未解除假设中自由出现。
\(P(x)\vdash P(x)\) 因而不允许直接推出 \(P(x)\vdash\forall xP(x)\)。

存在引入由 \(\varphi[t/x]\) 得到 \(\exists x\varphi\)。存在消去先取一个新变量
\(y\) 表示见证，以 \(\varphi[y/x]\) 为临时假设推出 \(\psi\)，再解除该假设。
\(y\) 不得在其他未解除假设、结论 \(\psi\) 或原存在公式中自由出现；代入仍须无捕获。
等号规则包括 \(t=t\)，以及在无捕获代入的公式中以相等项替换。

\section{两个不同的元定理}
\begin{theorem}{Soundness}
对上述经典演算，若 \(\Gamma\vdash\varphi\)，则 \(\Gamma\models\varphi\)。
\end{theorem}
\begin{proof}
对有限推导树归纳。假设在使背景假设全为真的解释中为真；联结词规则由真值规则验证。
全称引入时改变 \(x\) 的赋值不会改变背景假设的真值，这正是其附加条件的作用。
存在消去把新变量赋为一个见证，其附加条件保证背景与结论不依赖这个变量。
等号替换保持对象相同；反证规则使用经典二值语义。故根结论在每个这样的解释中为真。
\end{proof}

\begin{theorem}{Completeness of first-order logic}
对经典一阶逻辑的完备演算，\(\Gamma\models\varphi\) 蕴含 \(\Gamma\vdash\varphi\)。
\end{theorem}
\begin{qlnote}{证明路线}
对句子理论，设 \(\Gamma\nvdash\varphi\)。先证明 \(\Gamma\cup\{\neg\varphi\}\)
一致，再扩充语言加入存在见证，扩展为完备且一致的 Henkin 理论，最后以闭项的等价类
构造模型并证明 truth lemma。这样得到使 \(\Gamma\) 真而 \(\varphi\) 假的模型。
这里记录的是构造路线；见证扩充、一致性保留和 truth lemma 的完整证明须随原书逐步补足。
\end{qlnote}

\begin{theorem}{Compactness}
一阶句子集合 \(\Gamma\) 可满足，当且仅当其每个有限子集可满足。
\end{theorem}
\begin{proof}
一个模型当然满足每个有限子集。反之，若 \(\Gamma\) 不可满足，则
\(\Gamma\models\bot\)。由完备性有 \(\Gamma\vdash\bot\)，而这个有限推导只使用
有限多个来自 \(\Gamma\) 的假设。可靠性说明这些假设组成的有限子集不可满足，矛盾。
\end{proof}

逻辑完备性把语义后承与形式可推导性对应起来；它不会自动证明某个指定理论一致，
也不表示 ZFC 能决定其语言中的每个句子。不可判定性和不完备性留到 \emph{Logic}
Chapter Three 单独学习。
